\documentclass[10pt,a4paper]{article}
\usepackage[margin=24mm]{geometry}
\usepackage{amsmath,amssymb,booktabs}
\usepackage[T1]{fontenc}
\usepackage{lmodern}
\usepackage[hidelinks]{hyperref}
\setlength{\parindent}{0pt}
\setlength{\parskip}{4pt}
\newcommand{\Q}{\mathbb Q}
\newcommand{\R}{\mathbb R}
\newcommand{\C}{\mathbb C}
\newcommand{\Nm}{\operatorname{N}}
\begin{document}
\begin{center}
{\Large\bfseries Problem 00000002101: an arithmetic obstruction\\
to the printed square-root/unit conjunction}\par
\smallskip
\texttt{lizaixi01}\quad -- \quad 24 October 2026
\end{center}

\section*{Precise scope of the refutation}
The conjecture says that the Stokes constants are units of $\Q(\sqrt5)$
and prints the explicit value $\varphi^{1/2}$. We use the usual
golden-ratio convention
\[
 \varphi=\frac{1+\sqrt5}{2}.
\]
The source supplies no other numerical definition of $\varphi$.
We prove that \emph{every} complex square root of this number lies outside
the actual embedded quadratic field $\Q(\sqrt5)$. Consequently no scalar
can satisfy both the printed explicit-value and field-unit conditions.

This is an arithmetic obstruction to a necessary conjunction. We do not
assign a value to a Stokes constant, compute any actual Stokes data,
or disprove the field-unit assertion considered by itself. No assertion
against general DT--RH theory or an unstated alternative meaning of
$\varphi$ is made. Bridgeland's actual Kronecker BPS data and RH jump
conditions are given in Sections~2.9 and~3.1, printed pages~9--10 of
\cite{Bridgeland}; there the Stokes factors are Poisson automorphisms of
a torus. That reference provides context, not a proof or definition of
the conjectured numerical value.

\section*{A genuine quadratic field and its embeddings}
Construct
\[
 F=\Q[\omega],\qquad \omega^2=5,
\]
using the native \texttt{QuadraticAlgebra} with parameters $5,0$.
There is no rational $r$ with $r^2=5$: otherwise comparison with the two
real roots $\pm\sqrt5$ would contradict the irrationality of $\sqrt5$.
This proved fact supplies the actual native \texttt{Field} instance.
The genuine basis $1,\omega$ gives $[F:\Q]=2$.

The real embedding and actual conjugation are
\[
 \iota(a+b\omega)=a+b\sqrt5,
 \qquad \sigma(a+b\omega)=a-b\omega.
\]
Conjugation is an algebra automorphism, and $\iota$ is an injective algebra
homomorphism. Its range equals the native intermediate field
$\Q(\sqrt5)\subset\R$, defined as the field adjoined by $\sqrt5$.
Indeed the range contains $\sqrt5=\iota(\omega)$; conversely every
$a+b\sqrt5$ belongs to the adjoined field. These inclusions give an
actual algebra equivalence between $F$ and that full field, not just a
coordinate-vector identification. Composing with $\R\hookrightarrow\C$
gives the corresponding complex embedding; its range is likewise
exactly the actual complex intermediate field adjoined by $\sqrt5$.

\section*{The native field norm}
For $z=a+b\omega$, multiplication by $z$ on the genuine basis has matrix
\[
 M_z=\begin{pmatrix}a&5b\\b&a\end{pmatrix}.
\]
The native \texttt{Algebra.norm}, defined by this endomorphism's
determinant, is therefore
\[
 \Nm_{F/\Q}(a+b\omega)=a^2-5b^2,
 \qquad z\,\sigma(z)=\Nm_{F/\Q}(z)\cdot1.
\]
The proof explicitly identifies the quadratic-algebra norm with this
native determinant norm. It also transports the native norm through
the actual real and complex field equivalences, so the calculations
hold in both genuine adjoined fields.

Let $g=(1+\omega)/2$, concretely the field element with coordinates
$(1/2,1/2)$. Its real image is $\varphi$, its conjugate's real image is
$(1-\sqrt5)/2$, and the computed native norm is
\[
 \Nm_{F/\Q}(g)=\frac14-5\frac14=-1.
\]
The element $g$ itself is a genuine field unit, since its real image is
positive. The obstruction concerns its square roots.

\section*{All square roots are excluded}
If $x\in F$ satisfied $x^2=g$, multiplicativity of the native norm would
give
\[
  \Nm_{F/\Q}(x)^2=\Nm_{F/\Q}(x^2)=\Nm_{F/\Q}(g)=-1.
\]
The left side is a square of a rational number and is nonnegative, a
contradiction. Thus $g$ has no square root in $F$.

More generally, let $f:F\hookrightarrow K$ be \emph{any} field embedding
and suppose $c^2=f(g)$. If $c$ belonged to the range of $f$, write
$c=f(x)$. Injectivity gives $x^2=g$, which was just excluded. Apply this
to the actual real and complex embeddings. Every $c\in\C$ with
$c^2=\varphi$ therefore lies outside $\Q(\sqrt5)\subset\C$ and cannot
be the image of a unit of that actual intermediate field.

In particular the actual principal real value $s=\sqrt\varphi$ is positive
and satisfies $s^2=\varphi$. The real and complex square-root equations
have exactly the branches $s$ and $-s$; the argument excludes both.
It applies pointwise to every possible family of scalar data, so there
is no need to construct or assign a replacement Stokes object to obtain
the contradiction. Only the printed explicit-value/field-unit conjunction
under the stated golden-ratio convention is refuted.

\section*{Formalization and reproducibility}
The complete proof is in \texttt{lean/Main.lean}, namespace
\texttt{TLMC2101}.\par
Its main theorem is \texttt{complete\_counterexample}.
Principal bridges include:
\begin{center}
\begin{tabular}{ll}
\toprule
Bridge & Declaration\\
\midrule
Actual quadratic-field dimension & \texttt{quadratic\_degree}\\
True real adjoin-field identification & \texttt{realEmbedding\_range}\\
Native determinant norm identity & \texttt{nativeNorm\_eq\_quadraticNorm}\\
Norm in the actual complex field & \texttt{actual\_complex\_field\_golden\_norm}\\
All embedding images exclude roots & \texttt{no\_embedded\_square\_root}\\
Both complex branches & \texttt{complex\_square\_root\_branches}\\
Actual field-unit conjunction fails & \texttt{no\_unit\_square\_root}\\
\bottomrule
\end{tabular}
\end{center}
The project pins Lean~4.33.0 and Mathlib commit\par
\texttt{db584cd6d46c92f209a44c0f1c829460d327499d}.\par
Run \texttt{lean --version} separately; \texttt{lean/verify.sh} then runs
the root build, warning-as-error audit and official full fresh-kernel
replay. The audit enumerates every owned declaration, including generated
names, and runs ordinary \texttt{\#print axioms} on each.
Only \texttt{propext}, \texttt{Classical.choice} and \texttt{Quot.sound}
are allowed foundational axioms.

\begin{thebibliography}{9}
\bibitem{Bridgeland}
T. Bridgeland, \emph{Riemann--Hilbert problems from Donaldson--Thomas
theory}, arXiv:1611.03697v5 (24 June 2020), Sections~2.9 and~3.1,
printed pp.~9--10. \url{https://arxiv.org/abs/1611.03697v5}.
\end{thebibliography}
\end{document}
